\documentclass[a4paper,11pt,fleqn]{article}%fleqn pour aligner les equations à gauche
\usepackage{../../Styles/Exercice}

\newcommand{\chnum}{Chapitre 1}
\newcommand{\chtitle}{Réactions acide-base}
\newcommand{\dtype}{Exercices}

\begin{document}
\maketitle

\begin{multicols}{2}
\section{Une base de Bronsted est une espèce chimique susceptible de :}
\begin{enumerate}
    \item créer un ion hydrogène
    \item céder un ion hydrogène
    \item capter un ion hydrogène
\end{enumerate}

\section{Le schéma de Lewis de l'acide éthanoïque est :}
\begin{enumerate}
    \item \chemfig[atom sep = 2em]{H-C(-[2]H)(-[6]H)-C(=[1]\charge{0=\|,90=\|}{O})(-[7]\charge{45=\|,225=\|,315=\|,45:5pt=$\ominus$}{O})}
    \item \chemfig[atom sep = 2em]{H-C(-[2]H)(-[6]H)-C(=[1]O)(-[7]O-H)}
    \item \chemfig[atom sep = 2em]{H-C(-[2]H)(-[6]H)-C(=[1]\charge{0=\|,90=\|}{O})(-[7]\charge{180=\|,270=\|}{O}-H)}
\end{enumerate}

\section{Les espèces chimiques \ce{H2O}, \ce{H3O+} et \ce{HO-} forment des couples acide-base :}
\begin{enumerate}
    \item \ce{H3O+}/\ce{H2O} et \ce{H2O}/\ce{HO-}
    \item \ce{H3O+}/\ce{HO-} et \ce{H2O}/\ce{H2O}
    \item \ce{H2O}/\ce{H3O+} et \ce{HO-}/\ce{H2O}
\end{enumerate}

\section{L'espèce amphotère citée dans l'exercice précédent est :}
\begin{tabular}{m{.25\linewidth}m{.25\linewidth}m{.25\linewidth}}
    a) \ce{H3O+} & b) \ce{H2O} & c) \ce{HO-}\\
\end{tabular}

\section{Parmi les équations de réaction ci-dessous, celle qui correspond à une réaction acide-base est :}
\begin{enumerate}
    \item \ce{HCO2H(aq) + NH4+(aq) -> HCO2-(aq) + NH3(aq)}
    \item \ce{HCO2H(aq) + HCO2-(aq) -> HCO2-(aq) + HCO2H(aq)}
    \item \ce{HCO2H(aq) + NH3(aq) -> HCO2-(aq) + NH4+(aq)}
\end{enumerate}

\section{Une solution aqueuse a un pH de \num{8,0}. La concentration en ion oxonium, [\ce{H3O+}], en \unit{\mole\per\liter} est égale à :}
\begin{tabular}{m{.25\linewidth}m{.25\linewidth}m{.25\linewidth}}
    a) \num{1,0E-8} & b) \num{8,0E-1} & c) \num{1,0E8}\\
\end{tabular}


\section{La concentration en ions oxonium d'une solution aqueuse est [\ce{H3O+}] = \SI{1,0E-3}{\mole\per\liter}. Son pH est :}
\begin{tabular}{m{.25\linewidth}m{.25\linewidth}m{.25\linewidth}}
    a) \num{1,0} & b) \num{3,0} & c) \num{4,0}\\
\end{tabular}

\section{Lorsqu'un acide dest ajouté de façon modérée à une solution tampon, le pH de cette solution :}
\begin{enumerate}
    \item augmente très peu
    \item diminue très peu
    \item diminue beaucoup
\end{enumerate}

\section{Correspondances}
Recopier puis compléter le tableau ci-dessous sans utiliser de calculatrice.\\
\begin{tabular}{|m{.2\linewidth}|m{.15\linewidth}|m{.15\linewidth}|m{.16\linewidth}|m{.15\linewidth}|}
    \hline
    \textbf{[\ce{H3O+}] en \unit{\mole\per\liter}} & \small \num{1,0E-4} & & \small\num{1,0E-10} & \\
    \hline
    \textbf{pH} & & \small\num{5,0}& & \small\num{10,1}\\
    \hline
\end{tabular}
\textbf{Donnée :} $10^{-10,1}$ = \num{7,9E-11}

\section{Les fourmis}
Pour défendre leur nid, des fourmis aspergent d'acide méthanoïque leurs ennemis (essentiellement constitués d'eau)\\
\textbf{Données :} Couples acide-base : \ce{HCO2H}/\ce{HCO2-} et \ce{H3O+}/\ce{H2O}.
\begin{enumerate}
    \item Établie l'équation de la réaction acide-base à l'origine des brûlures. Identifier le transfert d'ion hydrogène à l'aide de schémas de Lewis adaptés.
    \item Préparer une communication orale sans note afin d'expliquer pourquoi n acide carboxylique est qualifié d'\textit{acide}. La démonstration doit s'appuyer sur du vocabulaire approprié et l'utilisation raisonnée d'outils de représentation des espèces chimiques.
\end{enumerate}

\section{Correction du pH de l'eau des piscines}
Le pH des piscines doit être compris entre 7,4 et 7,6. Pour le réguler, deux produits d'entretien sont disponibles :
\begin{itemize}
    \item produit A : hydrogénosulfate de sodium \ce{NaHSO4(s)}
    \item produit B , carbonate de calcium \ce{CaCO3(s)}
\end{itemize}
\textbf{Donnée :} couples acide-base : \ce{HSO4-}/\ce{SO4^{2-}} ; \ce{HCO3-}/\ce{CO3^{2-}} ; \ce{H3O+}/\ce{H2O} et \ce{H2O}/\ce{HO-}\\
Attrribuer à chaque produit son nim commercial "pH plus" ou "pH moins". Justifier en établissant les équations de dissolution puis les équations de réaction acide-base qui ont lieu lors de l'ajout de chaque procuit d'entretien dans l'eau d'une piscine.

\section{Calculer un pH}
L'addition de probure d'hydrogène gazeux \ce{HBr(g)} dnas de l'eau conduit à la formation d'acide bromhydrique. L'équation de la réaction acide-base qui a lieu est : 
$$ \ce{HBr(g) + H2O(l) -> Br-(aq) + H3O+(aq)}$$
Une solution aqueuse est préparée en faisant barboter du bromure d'hydrogène en quantité $n=\SI{0,39}{\milli\mole}$ afin d'obtenir une solution de volume $V=\SI{300}{mL}$.\\
Exprimer puis calculer le pH de la solution ainsi obtenue.

\end{multicols}
\end{document}